
\subsection{Топологии и сходимости в пространстве операторов L(E1,E2).Норма оператора. Полнота нормированного пространства  L(E1,E2) (в случае, когда E2 — банахово)}

\begin{theorem}[Эквивалентные определения ограниченности]\label{equiv-bounded-oper}
	Пусть $A: E_1 \to E_2$ --- линейный оператор. Следующие условия эквивалентны:
	\begin{enumerate}
		\item $A$ --- ограниченный (переводит ограниченные множества в ограниченные);
		\item $\exists K \geqslant 0: \forall x \in E_1 \quad \|Ax\|_{E_2} \leqslant K \|x\|_{E_1}$;
		\item Образ единичного шара $B = \{x : \|x\| \leqslant 1\}$ ограничен.
	\end{enumerate}
\end{theorem}

\begin{definition}[Норма оператора]
	\textbf{Нормой} линейного ограниченного оператора $A$ называется точная нижняя грань констант, ограничивающих оператор:
	\[ \|A\| = \inf \{ K : \forall x \in E_1, \|Ax\| \leqslant K \|x\| \} \]
\end{definition}

\begin{idea}
	Норма оператора $\|A\|$ — это «максимальный коэффициент растяжения». Это число показывает, во сколько раз оператор может увеличить длину вектора в худшем случае.
\end{idea}


\begin{proposition}[Вычисление нормы]
	Пусть $A$ --- линейный ограниченный оператор. Тогда:
	\[ \|A\| = \sup_{\|x\| \leqslant 1} \|Ax\| = \sup_{\|x\| = 1} \|Ax\| = \sup_{x \neq 0} \frac{\|Ax\|}{\|x\|} \]
\end{proposition}

\begin{definition}
	$\mathcal{L}(E_1, E_2)$ --- множество всех линейных ограниченных операторов из $E_1$ в $E_2$. Оно является линейным пространством с введенной выше нормой $\|A\|$.
\end{definition}

\begin{definition}[Сопряженное пространство]
	Пространство линейных непрерывных функционалов на $E$ называется \textbf{сопряженным (или двойственным)} пространством:
	\[ E^* = \mathcal{L}(E, \mathbb{K}) \]
\end{definition}

\begin{theorem}[О полноте пространства операторов]
	Пусть $E_1, E_2$ --- линейные нормированные пространства.
	\begin{enumerate}
		\item $\mathcal{L}(E_1, E_2)$ всегда является линейным нормированным пространством.
		\item Если $E_2$ --- \textbf{банахово} (полное), то $\mathcal{L}(E_1, E_2)$ тоже \textbf{банахово}.
	\end{enumerate}
\end{theorem}

\begin{remark}
	Полнота $E_1$ не требуется! Важна только полнота пространства, куда мы отображаем.
\end{remark}

\begin{proof}
	\textbf{1. Аксиомы нормы.}
	Линейность очевидна. Проверим неравенство треугольника:
	\[ \|A+B\| = \sup_{\|x\|=1} \|Ax + Bx\| \leqslant \sup_{\|x\|=1} (\|Ax\| + \|Bx\|) \leqslant \sup \|Ax\| + \sup \|Bx\| = \|A\| + \|B\|. \]
	
	\textbf{2. Полнота.}
	Пусть $E_2$ полно. Рассмотрим фундаментальную последовательность операторов $\{A_n\} \subset \mathcal{L}$.
	\[ \forall \varepsilon > 0 \exists N: \forall n, m \geqslant N \implies \|A_n - A_m\| < \varepsilon. \]
	
	\textbf{Шаг 1: Поточечная сходимость.}
	Зафиксируем произвольный $x \in E_1$.
	\[ \|A_n x - A_m x\| \leqslant \|A_n - A_m\| \cdot \|x\| \to 0. \]
	Значит, последовательность векторов $\{A_n x\}$ фундаментальна в $E_2$.
	Так как $E_2$ банахово, существует предел. Обозначим его $Ax := \lim_{n \to \infty} A_n x$.
	Оператор $A$ линеен (предел линейных отображений).
	
	\textbf{Шаг 2: Ограниченность предельного оператора $A$.}
	Фундаментальная последовательность $\{A_n\}$ ограничена по норме: $\exists C: \|A_n\| \leqslant C$.
	Тогда для любого $x$: $\|A_n x\| \leqslant C \|x\|$. Переходя к пределу: $\|Ax\| \leqslant C \|x\|$.
	Значит, $A \in \mathcal{L}(E_1, E_2)$.
	
	\textbf{Шаг 3: Сходимость по норме оператора ($\|A_n - A\| \to 0$).}
	В условии фундаментальности $\|A_n x - A_m x\| \leqslant \varepsilon \|x\|$ перейдем к пределу при $m \to \infty$:
	\[ \|A_n x - Ax\| \leqslant \varepsilon \|x\| \quad (\forall n \geqslant N, \forall x). \]
	Это означает, что $\sup_{\|x\|=1} \|(A_n - A)x\| \leqslant \varepsilon$, то есть $\|A_n - A\| \leqslant \varepsilon$.
	Следовательно, $A_n \to A$ в пространстве $\mathcal{L}$.
\end{proof}

\begin{corollary}
	Пространство $E^*$ всегда полное (так как $\mathbb{K} = \mathbb{R}$ или $\mathbb{C}$ полны).
\end{corollary}